\documentclass[11pt]{article}
\usepackage[letterpaper,margin=1.1in]{geometry}
\usepackage[T1]{fontenc}
\usepackage{lmodern}
\usepackage{microtype}
\usepackage{amsmath}
\usepackage{booktabs}
\usepackage{array}
\usepackage[round,authoryear]{natbib}
\usepackage{hyperref}
\hypersetup{colorlinks=true,linkcolor=black,citecolor=black,urlcolor=black,
  pdftitle={Earth Curvature and Refraction in the GoSurvey Traverse Editor},
  pdfauthor={Chet Jones}}
\usepackage{enumitem}
\setlength{\parindent}{0pt}
\setlength{\parskip}{0.55\baselineskip}
\newcolumntype{L}[1]{>{\raggedright\arraybackslash}p{#1}}

\title{\textbf{Earth Curvature and Refraction in the\\GoSurvey Traverse Editor}}
\author{Chet Jones \\ \small GoSurvey Project}
\date{October 2026}

\begin{document}
\maketitle
\vspace{-1.5em}
\begin{center}
\small REQ-826 \quad\textbar\quad GitHub Issue \#826 \quad\textbar\quad Pull Request \#839
\end{center}
\vspace{1em}

\begin{abstract}
\noindent GoSurvey's Traverse Editor reduced a leg's vertical angle to an
elevation difference ($\Delta Z$) by plain trigonometry alone. Over any leg
of meaningful length this is not what a level line actually measures: the
line of sight bends slightly toward the ground through atmospheric
refraction, and the reference surface itself curves away beneath it. This
note derives the standard combined correction, works it by hand for several
leg lengths and refraction conditions, and compares each hand result
against the value GoSurvey now computes. It also documents the automated
test suite that keeps the implementation honest going forward, and records
the resulting change in previously-computed traverse elevations.
\end{abstract}

\section{What problem this fixes}

A total station measures a vertical angle and a distance to a target, and
trigonometry converts that pair into a horizontal distance and an elevation
difference. Plain trigonometry implicitly assumes two things that are not
true over any real distance: that the line of sight travels in a straight
line through the atmosphere, and that ``level'' is a flat plane rather than
a surface curving with the earth. Both effects are well documented in
standard surveying references \citep{ghilaniwolf2015,bannister1998}, and
both act in the same direction on a normal (non-grazing) sight: curvature
makes a level target appear to drop away below the horizontal, while
atmospheric refraction bends the line of sight downward toward denser air
near the ground, partially canceling the curvature effect. Left
uncorrected, this is equivalent to always observing with a refraction
coefficient of $k \approx 1$ no matter what the atmosphere is actually
doing, an error that grows with the square of the distance and was reported
against GoSurvey directly by a user running real traverse legs
\citep{issue826}. This change adds the standard combined curvature and
refraction correction to the Traverse Editor's $\Delta Z$ computation.

\section{The governing formula}

The combined curvature and refraction correction to an elevation difference
observed over a horizontal distance $D$ is
\begin{equation}
C \;=\; (1-k)\,\frac{D^{2}}{2R},
\label{eq:cr}
\end{equation}
where $R$ is the mean radius of the earth and $k$ is the refraction
coefficient, the fraction of the curvature effect that atmospheric
refraction cancels out \citep{ghilaniwolf2015,bannister1998}. GoSurvey uses
$R = 6{,}371{,}000$~m, the mean earth radius, and a default
$k = 0.14$, the value standard references give for average atmospheric
conditions \citep{ghilaniwolf2015}. $C$ is strictly positive whenever
$k < 1$ (the normal case): the corrected elevation difference is always
\emph{higher} than the plain-trigonometric value, for a level or an
inclined sight in either direction, because curvature and refraction act on
the vertical reduction independently of which way the sight is pointed.

For each traverse leg, GoSurvey already reduces the field observation
(whichever of a direct horizontal distance or a slope distance plus
vertical angle was entered) to a horizontal distance $D$ and a plain-trig
elevation difference $\Delta Z_{\mathrm{trig}}$. This change adds one term
after that reduction, at the single point both input paths already share:
\begin{equation}
\Delta Z \;=\; \Delta Z_{\mathrm{trig}} \;+\; (1-k)\,\frac{D^{2}}{2R},
\label{eq:dz}
\end{equation}
applied only when a vertical angle was actually resolved for the leg; a leg
entered with no vertical observation keeps its legacy flat ($\Delta Z = 0$)
behavior exactly as before. The option defaults on, since the uncorrected
behavior is a systematic error rather than a deliberate choice, and can be
switched off to reproduce the exact pre-change numbers.

\emph{Implementation.} \texttt{src/traverse/TraverseCalc.cpp}, function
\texttt{ComputeTraverse}, in the per-leg loop immediately after the
horizontal distance is resolved and validated, before it is used for
anything else (coordinate increments, closure perimeter, or the
least-squares adjustment's inputs). \texttt{src/traverse/TraverseLeastSquares.cpp}
performs no vertical reduction of its own, so this single point covers
every elevation the Traverse Editor produces.

\section{Worked example A: the textbook one-kilometer sight}
\label{sec:exa}

A $1000.0000$~m leg with a $+2^{\circ}$ elevation angle, at the default
$k = 0.14$:

\begin{center}
\begin{tabular}{@{}llr@{}}
\toprule
Quantity & Formula & Value \\
\midrule
Horizontal distance $D$ & measured & $1000.0000$ m \\
Plain-trig $\Delta Z_{\mathrm{trig}}$ & $D\tan(2^{\circ})$ & $34.9207$ m \\
Curvature + refraction $C$ & $(1-0.14)D^{2}/(2R)$ & $0.0674933$ m \\
Corrected $\Delta Z$ & $\Delta Z_{\mathrm{trig}} + C$ & $34.9882$ m \\
\bottomrule
\end{tabular}
\end{center}

This is the textbook anchor value: $C = 0.0675\,K^{2}$ for $K$ in kilometers
at $k=0.14$ \citep{ghilaniwolf2015}, and $0.0675 \times 1^2 = 0.0675$~m
agrees with the hand-computed $0.0674933$~m to the fourth decimal (the
small residual is simply the published constant's own rounding). GoSurvey's
test suite asserts this exact pair of values:
\texttt{tests/TraverseTests.cpp}, \emph{``Curvature+refraction: 1 km leg
gains 0.0675 m at k=0.14''}, which isolates the correction by differencing a
corrected run against the same leg with the option off, so the test is
immune to any unrelated change in the plain-trig term.

\section{Worked example B: refraction coefficient and leg length both matter}
\label{sec:exb}

Because $C$ depends only on $D$ and $k$, not on the vertical angle itself,
it is straightforward to check the formula's two free parameters
independently against hand arithmetic.

\begin{center}
\begin{tabular}{@{}L{0.30\textwidth}L{0.34\textwidth}rr@{}}
\toprule
Case & Formula & Hand value & GoSurvey \\
\midrule
$D=1000$~m, $k=0$ (pure curvature, no refraction) &
  $D^{2}/(2R)$ & $0.0784804$ m & $0.078480$ m \\
$D=1000$~m, $k=0.20$ (hazier air, more bending) &
  $0.80\,D^{2}/(2R)$ & $0.0627843$ m & $0.062784$ m \\
$D=300$~m, $k=0.14$ (a short layout leg) &
  $0.86\,D^{2}/(2R)$ & $0.0060744$ m & $0.0060744$ m \\
\bottomrule
\end{tabular}
\end{center}

The $k=0$ row is the full, uncancelled curvature drop, the largest the
correction can be for a given distance under this formula; the $k=0.20$ row
shows the correction shrinking as $k$ rises toward $1$, exactly as
Equation~\ref{eq:cr} predicts; and the $300$~m row shows the term is
already under a centimeter at ordinary traverse-leg lengths, which is why
the uncorrected error was easy to miss on short legs and only became
obvious on the longer legs the original bug report described
\citep{issue826}. Each row is covered by its own case in
\texttt{tests/TraverseTests.cpp} under the \texttt{[curvref][issue-826]}
tag.

\subsection*{Sign check: downhill sights}

Equation~\ref{eq:dz} adds $C$ unconditionally, so a downhill
($-2^{\circ}$) sight over the same $1000$~m gets the identical
$+0.0674933$~m correction as the uphill case in
Section~\ref{sec:exa} --- the corrected $\Delta Z$ is less negative than
the plain-trig value, not more. This matches the physical picture: a level
line of sight over a curved earth drops below the tangent plane regardless
of which way the instrument is tilted, so the curvature term must raise
the elevation difference in both directions. GoSurvey's test
\emph{``Curvature+refraction: downhill sights get a positive correction''}
asserts both the magnitude and the direction of this effect.

\subsection*{The slope-distance path}

A leg entered as a $1000$~m slope distance at a zenith angle of
$88^{\circ}$ reduces to a horizontal distance of
$D = 1000\sin(88^{\circ}) = 999.3908$~m before the correction term is
applied, giving
$C = 0.86 \times 999.3908^{2}/(2\times 6{,}371{,}000) = 0.0674111$~m, which
GoSurvey's \emph{``Curvature+refraction: slope-distance/zenith path
matches''} test confirms to the same precision. This matters because
GoSurvey has two independent input paths for a leg (an entered horizontal
distance, or a slope distance reduced through a separate function); the
test demonstrates both paths reach the same correction code rather than
only one of them being fixed.

\section{The default (no vertical angle) path is unchanged}

A leg entered with only a horizontal distance and no vertical
observation --- there being nothing to reduce --- keeps
$\Delta Z = 0.0000000$~m exactly, bit-for-bit identical to the pre-change
behavior, confirmed by \emph{``Curvature+refraction: legs with no vertical
angle are untouched''}. This guards against a plausible implementation
mistake: applying the curvature term to every leg's horizontal distance
regardless of whether a vertical angle exists, which would silently invent
elevation data for legs that never had any.

\section{What changes for existing users}

Because the option defaults on, elevations computed by the Traverse Editor
for any leg that used a vertical angle will differ from earlier GoSurvey
versions by exactly the $C$ term above: roughly $+6.75$~cm per kilometer of
leg length at the default $k=0.14$, falling to about $+6.1$~mm at
$300$~m and $+0.7$~mm at $100$~m. This is not a new bug; it corrects an
existing one. Unchecking ``Curv + Refr'' in the Traverse Editor's options
row restores the exact previous numbers, verified by the no-op test in
Section~5 applied with the toggle left at its non-default (off) state.

\section{Automated test summary}

\begin{center}
\small
\begin{tabular}{@{}L{0.27\textwidth}L{0.37\textwidth}L{0.28\textwidth}@{}}
\toprule
File & Test case & What it proves \\
\midrule
\texttt{TraverseTests.cpp} &
  1 km leg gains 0.0675 m at k=0.14 &
  the textbook anchor value, isolated from the plain-trig term \\[0.3em]
\texttt{TraverseTests.cpp} &
  k=0 gives pure curvature 0.0785 m/km &
  the uncancelled upper bound of the correction \\[0.3em]
\texttt{TraverseTests.cpp} &
  custom k=0.20 scales the correction &
  $k$ is a live user input, not a baked-in constant \\[0.3em]
\texttt{TraverseTests.cpp} &
  300 m leg corrects 6.1 mm &
  the term stays correct, and small, at ordinary leg lengths \\[0.3em]
\texttt{TraverseTests.cpp} &
  downhill sights get a positive correction &
  the sign is direction-independent, as the physics requires \\[0.3em]
\texttt{TraverseTests.cpp} &
  slope-distance/zenith path matches &
  both leg-entry paths reach the one correction point \\[0.3em]
\texttt{TraverseTests.cpp} &
  legs with no vertical angle are untouched &
  the correction cannot invent elevation data from nothing \\
\bottomrule
\end{tabular}
\end{center}

All seven cases, plus the eight pre-existing traverse test cases, pass
under the \texttt{GoSurveyTests} executable (Catch2, MSVC release build),
15 test cases and 97 assertions with no regressions. The source is
PR~\#839, branch \texttt{fix-826-curv-refr}, requirement entry pending in
\texttt{spec/requirements.md}.

\begin{thebibliography}{9}

\bibitem[Bannister et al., 1998]{bannister1998}
Bannister, A., Raymond, S., and Baker, R. (1998). \emph{Surveying}, 7th
ed. Harlow, UK: Pearson/Prentice Hall.

\bibitem[Ghilani and Wolf, 2015]{ghilaniwolf2015}
Ghilani, C.~D., and Wolf, P.~R. (2015). \emph{Elementary Surveying: An
Introduction to Geomatics}, 14th ed. Boston: Pearson.

\bibitem[GoSurvey Issue 826, 2026]{issue826}
GoSurvey Project (2026). \emph{GitHub Issue \#826: Traverse editor: add
earth curvature + refraction correction to vertical angle reduction}.
\url{https://github.com/chetjones003/GoSurvey/issues/826}

\end{thebibliography}

\end{document}
